%% @file paper53_differentialgeometrie_de.tex
%% @brief Differentialgeometrie: Mannigfaltigkeiten, Krümmung und das Zusammenspiel von Geometrie und Analysis
%% @author Michael Fuhrmann
%% @date 2026-03-12
%% @project Specialist Mathematics Project, Build 139

\documentclass[12pt,a4paper]{amsart}
\usepackage{amsmath,amssymb,amsthm}
\usepackage{mathtools}
\usepackage[utf8]{inputenc}
\usepackage[T1]{fontenc}
\usepackage[ngerman]{babel}
\usepackage{hyperref}
\usepackage{enumitem,booktabs,array}
\usepackage{geometry}
\geometry{margin=2.5cm}

\newtheorem{theorem}{Satz}[section]
\newtheorem{lemma}[theorem]{Lemma}
\newtheorem{korollar}[theorem]{Korollar}
\newtheorem{proposition}[theorem]{Proposition}
\newtheorem{vermutung}[theorem]{Vermutung}
\theoremstyle{definition}
\newtheorem{definition}[theorem]{Definition}
\newtheorem{beispiel}[theorem]{Beispiel}
\newtheorem{bemerkung}[theorem]{Bemerkung}

\author{Michael Fuhrmann}
\address{Specialist Mathematics Project, Build 139}
\email{specialist-maths@localhost}
\date{\today}

\begin{document}
\title{Differentialgeometrie: Mannigfaltigkeiten, Kr\"ummung und das\\
Zusammenspiel von Geometrie und Analysis}

\begin{abstract}
Die Differentialgeometrie ist das mathematische Teilgebiet, das glatte
Mannigfaltigkeiten und ihre geometrischen Eigenschaften mithilfe der
Differentialrechnung und Tensoranalysis untersucht. Dieser Artikel gibt
eine strenge und zugleich zugängliche Einführung in die Kernkonzepte
der modernen Differentialgeometrie: glatte Mannigfaltigkeiten mit Karten
und Atlanten, Riemannsche Metriken zur Messung von Längen und Winkeln,
Zusammenhänge und kovariante Ableitungen --- aufgipfelnd im
Levi-Civita-Zusammenhang (als eindeutig bewiesen durch Metrik-Kompatibilität
und Torsionsfreiheit) --- sowie den Riemann-Krümmungstensor mit seinen
abgeleiteten Größen: Ricci-Tensor, Skalarkrümmung und Gaußsche Krümmung.
Der klassische Gauß-Bonnet-Satz, der lokale Krümmung mit dem globalen
topologischen Invarianten (Euler-Charakteristik) verbindet, wird für
geschlossene Flächen bewiesen; seine weitreichende Verallgemeinerung, der
Chern-Gauß-Bonnet-Satz für gerade-dimensionale Mannigfaltigkeiten, wird
ebenfalls vollständig dargestellt. Es werden Geodäten über die Geodätengleichung
und den Satz von Hopf-Rinow (bewiesen), Jacobi-Felder sowie die zweite
Fundamentalform für Untermannigfaltigkeiten behandelt. Lie-Gruppen als
geometrische Objekte mit linksinvarianten Metriken bieten eine algebraische
Perspektive. Der Artikel schließt mit physikalischen Anwendungen in der
Allgemeinen Relativitätstheorie, wo die Einsteinschen Feldgleichungen
die Raumzeitkrümmung durch Materie und Energie beschreiben. Das begleitende
Python-Modul \texttt{differential\_geometry.py} implementiert die hier
besprochenen Algorithmen und liefert numerische Verifikation der theoretischen
Ergebnisse.
\end{abstract}

\maketitle
\tableofcontents

% ============================================================
\section{Einleitung}
% ============================================================

Die Differentialgeometrie nimmt einen zentralen Platz in der modernen Mathematik
ein und bildet die natürliche Sprache sowohl für rein mathematische Strukturen
als auch für die Grundgesetze der Physik. Ihr Kern liegt in einer scheinbar
einfachen Frage: Wie misst und vergleicht man geometrische Größen auf gekrümmten
Räumen, die sich nicht kanonisch in einen flachen Umgebungsraum einbetten lassen?

Die Ursprünge des Fachgebiets reichen zurück zu Carl Friedrich Gauß (1777--1855),
dessen \emph{Disquisitiones generales circa superficies curvas} (1827) die
revolutionäre Idee einführte, dass die \emph{innere Geometrie} einer Fläche---
die Geometrie, wie sie ein Wesen erlebt, das auf der Fläche lebt---vollständig
durch die erste Fundamentalform bestimmt ist, ohne Bezug auf einen Umgebungsraum.
Dies ist der Inhalt von Gauß' \emph{Theorema Egregium}: Die Gaußsche Krümmung $K$
ist eine innere Invariante.

Bernhard Riemann (1826--1866) verallgemeinerte Gauß' Ideen in seinem berühmten
Habilitationsvortrag von 1854 \emph{Über die Hypothesen welche der Geometrie
zu Grunde liegen} auf Mannigfaltigkeiten beliebiger Dimension. Riemann führte
den metrischen Tensor $g_{ij}$ ein, der die Messung von Längen, Winkeln und
Volumina auf glatten Mannigfaltigkeiten ermöglicht. Seine Vision legte den Grundstein für:
\begin{itemize}
  \item Die intrinsische Theorie glatter Mannigfaltigkeiten (Whitney, 1936),
  \item Die Riemannsche Geometrie mit metrischen Tensoren,
  \item Den kovarianten Kalkül von Gregorio Ricci-Curbastro und Tullio
        Levi-Civita (1900--1917),
  \item Die Allgemeine Relativitätstheorie (Einstein, 1915), wo Raumzeit
        eine 4-dimensionale Lorentz-Mannigfaltigkeit ist.
\end{itemize}

Das Wechselspiel zwischen Geometrie und Analysis ist allgegenwärtig. Differentialformen
verbinden Krümmung und Topologie über die de-Rham-Kohomologie. Das Spektrum des
Laplace-Beltrami-Operators kodiert geometrische Information. Minimalflächen ---
Lösungen von Variationsproblemen --- erscheinen in Seifenblasenexperimenten
ebenso wie in der modernen Stringtheorie.

\subsection{Aufbau dieses Artikels}

Abschnitt~\ref{sec:mannigfaltigkeiten} führt glatte Mannigfaltigkeiten rigoros
über Karten und Atlanten ein. Abschnitt~\ref{sec:metriken} entwickelt Riemannsche
Metriken, Längen und die Exponentialabbildung. Abschnitt~\ref{sec:zusammenhaenge}
behandelt Zusammenhänge und kovariante Ableitungen, mit der Eindeutigkeit des
Levi-Civita-Zusammenhangs als Höhepunkt. Abschnitt~\ref{sec:kruemmung} analysiert
die vollständige Hierarchie von Krümmungstensoren.
Abschnitt~\ref{sec:gauss-bonnet} präsentiert den Gauß-Bonnet-Satz und Cherns
Verallgemeinerung. Abschnitt~\ref{sec:geodaeten} untersucht Geodäten,
Vollständigkeit und Jacobi-Felder. Abschnitt~\ref{sec:untermannigfaltigkeiten}
führt die zweite Fundamentalform und Minimalflächen ein.
Abschnitt~\ref{sec:liegruppen} verbindet Lie-Gruppen mit der Geometrie.
Abschnitt~\ref{sec:anwendungen} beschreibt Anwendungen in der Allgemeinen
Relativitätstheorie.

% ============================================================
\section{Glatte Mannigfaltigkeiten}
\label{sec:mannigfaltigkeiten}
% ============================================================

\begin{definition}[Topologische Mannigfaltigkeit]
Eine \emph{topologische Mannigfaltigkeit} $M$ der Dimension $n$ ist ein
zweitabzählbarer Hausdorff'scher topologischer Raum, sodass jeder Punkt
$p \in M$ eine offene Umgebung $U$ besitzt, die homöomorph zu einer
offenen Teilmenge von $\mathbb{R}^n$ ist.
\end{definition}

\begin{definition}[Karte und Atlas]
Eine \emph{Karte} (oder \emph{Koordinatenkarte}) auf $M$ ist ein Paar
$(U, \varphi)$, wobei $U \subseteq M$ offen und $\varphi: U \to V \subseteq \mathbb{R}^n$
ein Homöomorphismus ist. Ein \emph{glatter Atlas} ist eine Kollektion
$\mathcal{A} = \{(U_\alpha, \varphi_\alpha)\}$ von Karten mit:
\begin{enumerate}[label=(\roman*)]
  \item $\bigcup_\alpha U_\alpha = M$ (Überdeckungsbedingung),
  \item Für je zwei Karten $(U_\alpha, \varphi_\alpha)$ und
        $(U_\beta, \varphi_\beta)$ mit $U_\alpha \cap U_\beta \neq \emptyset$
        ist die \emph{Übergangsfunktion}
        \[
          \varphi_\beta \circ \varphi_\alpha^{-1}:
          \varphi_\alpha(U_\alpha \cap U_\beta) \to
          \varphi_\beta(U_\alpha \cap U_\beta)
        \]
        ein $C^\infty$-Diffeomorphismus.
\end{enumerate}
\end{definition}

\begin{definition}[Glatte Mannigfaltigkeit]
Eine \emph{glatte Mannigfaltigkeit} ist eine topologische Mannigfaltigkeit $M$,
ausgestattet mit einem maximalen glatten Atlas $\mathcal{A}$ (auch
\emph{glatte Struktur} genannt).
\end{definition}

\begin{beispiel}
\begin{enumerate}[label=(\alph*)]
  \item Die $n$-Sphäre $S^n = \{x \in \mathbb{R}^{n+1} : |x| = 1\}$ ist eine
        kompakte glatte $n$-Mannigfaltigkeit. Der Standardatlas besteht aus
        zwei stereographischen Projektionskarten.
  \item Der Torus $\mathbb{T}^2 = S^1 \times S^1$ ist eine kompakte 2-Mannigfaltigkeit
        mit Karten durch Winkelpaare $(\theta_1, \theta_2)$.
  \item Die allgemeine lineare Gruppe $\mathrm{GL}(n, \mathbb{R}) \subset M_n(\mathbb{R})$
        ist eine offene Teilmannigfaltigkeit von $\mathbb{R}^{n^2}$.
\end{enumerate}
\end{beispiel}

\subsection{Tangential- und Kotangentialraum}

\begin{definition}[Tangentialraum]
Sei $p \in M$. Eine \emph{Derivation} bei $p$ ist eine lineare Abbildung
$v: C^\infty(M) \to \mathbb{R}$, die die Leibniz-Regel erfüllt:
\[
  v(fg) = f(p)\,v(g) + g(p)\,v(f).
\]
Die Menge aller Derivationen bei $p$ ist der \emph{Tangentialraum} $T_pM$,
ein $n$-dimensionaler reeller Vektorraum. In einer Karte $(U, \varphi)$ mit
Koordinaten $(x^1, \ldots, x^n)$ sind die \emph{Koordinatenbasisvektoren}:
\[
  \left.\frac{\partial}{\partial x^i}\right|_p f
  = \left.\frac{\partial (f \circ \varphi^{-1})}{\partial r^i}\right|_{\varphi(p)},
  \quad i = 1, \ldots, n.
\]
\end{definition}

\begin{definition}[Kotangentialraum und Differentialformen]
Der \emph{Kotangentialraum} bei $p$ ist das Duale $T_p^*M = (T_pM)^*$.
Das Differential einer glatten Funktion $f: M \to \mathbb{R}$ ist
$(df)_p \in T_p^*M$, definiert durch $(df)_p(v) = v(f)$ für $v \in T_pM$.
Die Koordinatendifferentiale $\{dx^1|_p, \ldots, dx^n|_p\}$ bilden die
Dualbasis zu $\{\partial/\partial x^i|_p\}$.
\end{definition}

\begin{definition}[Tangentialbündel]
Das \emph{Tangentialbündel} ist $TM = \bigsqcup_{p \in M} T_pM$ mit der
natürlichen Projektion $\pi: TM \to M$. Ein \emph{glattes Vektorfeld}
auf $M$ ist ein glatter Schnitt $X: M \to TM$.
\end{definition}

\subsection{Parametrisierte Kurven und der Frenet-Serret-Rahmen}

Eine glatte Kurve in $\mathbb{R}^n$ ist eine glatte Abbildung
$\mathbf{r}: I \to \mathbb{R}^n$. Für eine reguläre Kurve ($|\mathbf{r}'(t)| > 0$)
besteht der \emph{Frenet-Serret-Rahmen} $\{T, N, B\}$ aus Tangenteneinheitsvektor,
Hauptnormalenvektor und Binormalvektor. Die \emph{Frenet-Serret-Gleichungen}
in $\mathbb{R}^3$ lauten:
\[
  T' = \kappa N, \quad N' = -\kappa T + \tau B, \quad B' = -\tau N,
\]
wobei $\kappa \geq 0$ die \emph{Krümmung} und $\tau \in \mathbb{R}$ die
\emph{Torsion} ist. Explizit gilt:
\[
  \kappa = \frac{|\mathbf{r}' \times \mathbf{r}''|}{|\mathbf{r}'|^3},
  \qquad
  \tau = \frac{(\mathbf{r}' \times \mathbf{r}'') \cdot \mathbf{r}'''}{|\mathbf{r}' \times \mathbf{r}''|^2}.
\]

\begin{beispiel}[Helix]
Die Helix $\mathbf{r}(t) = (\cos t, \sin t, t)$ hat
$\kappa = 1/2$ und $\tau = 1/2$ (beide konstant), was eine Kurve beschreibt,
die sich gleichmäßig um einen Zylinder windet.
\end{beispiel}

% ============================================================
\section{Riemannsche Metriken}
\label{sec:metriken}
% ============================================================

\begin{definition}[Riemannsche Metrik]
Eine \emph{Riemannsche Metrik} auf einer glatten Mannigfaltigkeit $M$ ist eine
glatte Zuordnung eines Skalarprodukts $g_p: T_pM \times T_pM \to \mathbb{R}$
an jeden Tangentialraum. In lokalen Koordinaten $(x^1, \ldots, x^n)$:
\[
  g = g_{ij}\,dx^i \otimes dx^j, \quad g_{ij} = g_{ji}, \quad (g_{ij}) \succ 0.
\]
Das Paar $(M, g)$ heißt \emph{Riemannsche Mannigfaltigkeit}.
\end{definition}

\begin{beispiel}[Klassische Metriken]
\begin{enumerate}[label=(\alph*)]
  \item \textbf{Euklidischer Raum}: $g = \delta_{ij}\,dx^i\,dx^j = dx^2 + dy^2 + dz^2$.
  \item \textbf{Runde Sphäre $S^2(r)$}: In sphärischen Koordinaten
        $(\theta, \phi)$: $g = r^2(d\theta^2 + \sin^2\!\theta\,d\phi^2)$.
  \item \textbf{Poincaré-Halbebene}: $\mathbb{H}^2 = \{(x,y): y>0\}$,
        $g = (dx^2 + dy^2)/y^2$. Konstante Krümmung $K = -1$.
  \item \textbf{Flacher Torus $\mathbb{T}^2$}: $g = d\theta_1^2 + d\theta_2^2$.
\end{enumerate}
\end{beispiel}

\subsection{Bogenlänge und die erste Fundamentalform}

Für eine parametrisierte Fläche $\mathbf{r}(u,v): U \to \mathbb{R}^3$ lautet
die \emph{erste Fundamentalform}:
\[
  \mathrm{I} = E\,du^2 + 2F\,du\,dv + G\,dv^2,
\]
mit $E = \mathbf{r}_u \cdot \mathbf{r}_u$, $F = \mathbf{r}_u \cdot \mathbf{r}_v$,
$G = \mathbf{r}_v \cdot \mathbf{r}_v$. Die \emph{Bogenlänge} einer Kurve
$\mathbf{r}(t)$ ist:
\[
  L = \int_{t_0}^{t_1} |\mathbf{r}'(t)|\,dt
    = \int_{t_0}^{t_1} \sqrt{g_{ij}\,\dot{x}^i\,\dot{x}^j}\,dt.
\]

\subsection{Die Exponentialabbildung}

\begin{definition}[Exponentialabbildung]
Sei $(M, g)$ eine Riemannsche Mannigfaltigkeit und $p \in M$. Für $v \in T_pM$
bezeichne $\gamma_v$ die eindeutige Geodäte mit $\gamma_v(0) = p$ und $\gamma_v'(0) = v$.
Die \emph{Exponentialabbildung} bei $p$ ist:
\[
  \exp_p: T_pM \supseteq \mathcal{D}_p \to M, \quad
  \exp_p(v) = \gamma_v(1).
\]
\end{definition}

\begin{proposition}
Für jeden Punkt $p \in M$ ist $\exp_p$ ein lokaler Diffeomorphismus von einer
Umgebung von $0 \in T_pM$ auf eine Umgebung von $p \in M$.
\end{proposition}

Normalkoordinaten um $p$, die über $\exp_p$ gewonnen werden, erfüllen
$g_{ij}(p) = \delta_{ij}$ und $\partial_k g_{ij}(p) = 0$, was grundlegend
für viele lokale Berechnungen ist.

% ============================================================
\section{Zusammenhänge und kovariante Ableitungen}
\label{sec:zusammenhaenge}
% ============================================================

Ein Zusammenhang (engl.\ \emph{connection}) liefert das Werkzeug, um Vektorfelder
entlang Kurven auf einer Mannigfaltigkeit zu differenzieren --- eine Verallgemeinerung
der Richtungsableitung von Vektorfeldern in $\mathbb{R}^n$.

\begin{definition}[Affiner Zusammenhang]
Ein \emph{affiner Zusammenhang} auf $M$ ist eine Abbildung
$\nabla: \mathfrak{X}(M) \times \mathfrak{X}(M) \to \mathfrak{X}(M)$,
$(X, Y) \mapsto \nabla_X Y$, die folgende Axiome erfüllt:
\begin{enumerate}[label=(\roman*)]
  \item $\nabla_{fX+gY} Z = f \nabla_X Z + g \nabla_Y Z$ \quad ($C^\infty$-Linearität in $X$),
  \item $\nabla_X(Y + Z) = \nabla_X Y + \nabla_X Z$,
  \item $\nabla_X(fY) = (Xf)\,Y + f\,\nabla_X Y$ (Leibniz-Regel).
\end{enumerate}
In lokalen Koordinaten gilt $\nabla_{\partial_i} \partial_j = \Gamma^k_{ij} \partial_k$,
wobei $\Gamma^k_{ij}$ die \emph{Christoffel-Symbole} sind.
\end{definition}

\begin{definition}[Metrik-Kompatibilität und Torsionsfreiheit]
Ein Zusammenhang $\nabla$ auf $(M, g)$ ist:
\begin{itemize}
  \item \emph{metrik-kompatibel}, falls $\nabla g = 0$, d.\,h.
        $X\bigl(g(Y,Z)\bigr) = g(\nabla_X Y, Z) + g(Y, \nabla_X Z)$,
  \item \emph{torsionsfrei} (symmetrisch), falls der Torsionstensor
        $T(X,Y) = \nabla_X Y - \nabla_Y X - [X,Y] = 0$ verschwindet.
\end{itemize}
\end{definition}

\begin{theorem}[Levi-Civita-Zusammenhang --- Fundamentalsatz der Riemannschen Geometrie]
\label{thm:levi-civita}
Sei $(M, g)$ eine Riemannsche Mannigfaltigkeit. Es existiert ein \textbf{eindeutiger}
affiner Zusammenhang $\nabla$ auf $M$, der sowohl metrik-kompatibel als auch
torsionsfrei ist. Dieser heißt \emph{Levi-Civita-Zusammenhang}.
\end{theorem}

\begin{proof}
\textbf{Eindeutigkeit.} Sei $\nabla$ metrik-kompatibel und torsionsfrei.
Zyklische Anwendung der Metrik-Kompatibilität liefert:
\begin{align*}
  Xg(Y,Z) &= g(\nabla_X Y, Z) + g(Y, \nabla_X Z), \\
  Yg(Z,X) &= g(\nabla_Y Z, X) + g(Z, \nabla_Y X), \\
  Zg(X,Y) &= g(\nabla_Z X, Y) + g(X, \nabla_Z Y).
\end{align*}
Addition der ersten beiden und Subtraktion der dritten, kombiniert mit
Torsionsfreiheit $\nabla_X Y - \nabla_Y X = [X,Y]$, ergibt die
\emph{Koszul-Formel}:
\begin{multline*}
  2\,g(\nabla_X Y, Z) = Xg(Y,Z) + Yg(X,Z) - Zg(X,Y) \\
    + g([X,Y], Z) - g([X,Z], Y) - g([Y,Z], X).
\end{multline*}
Da $g$ nicht-ausgeartet ist, ist $\nabla_X Y$ eindeutig bestimmt.

\textbf{Existenz.} Man definiert $\nabla_X Y$ durch die Koszul-Formel und
verifiziert direkt alle Verbindungsaxiome sowie Metrik-Kompatibilität und
Torsionsfreiheit. Die zugehörigen Christoffel-Symbole lauten:
\[
  \Gamma^k_{ij} = \tfrac{1}{2} g^{kl}
  \bigl(\partial_i g_{jl} + \partial_j g_{il} - \partial_l g_{ij}\bigr).
\]
\end{proof}

\subsection{Paralleltransport}

\begin{definition}[Paralleltransport]
Ein Vektorfeld $V$ entlang einer Kurve $\gamma: [a,b] \to M$ heißt
\emph{parallel} (längs $\gamma$), falls $\nabla_{\gamma'}V = 0$. In Koordinaten:
\[
  \frac{dV^k}{dt} + \Gamma^k_{ij}\,\frac{dx^i}{dt}\,V^j = 0.
\]
Der \emph{Paralleltransport} längs $\gamma$ ist der Isomorphismus
$P_{a \to t}: T_{\gamma(a)}M \to T_{\gamma(t)}M$, der jeden Tangentialvektor
bei $\gamma(a)$ in sein Paralleltranslatiertes bei $\gamma(t)$ überführt.
\end{definition}

\begin{bemerkung}
Paralleltransport auf einer gekrümmten Mannigfaltigkeit ist wegabhängig.
Die Holonomie einer geschlossenen Kurve auf $S^2$ ergibt eine Rotation gleich
dem eingeschlossenen Raumwinkel, direkt sichtbar als Präzession des
Foucault-Pendels auf der Erde.
\end{bemerkung}

% ============================================================
\section{Krümmung}
\label{sec:kruemmung}
% ============================================================

Krümmung misst das Scheitern der Vertauschbarkeit zweiter Ableitungen ---
die infinitesimale Version des wegabhängigen Paralleltransports.

\begin{definition}[Riemann-Krümmungstensor]
Der \emph{Riemann-Krümmungstensor} ist der $(1,3)$-Tensor:
\[
  R(X,Y)Z = \nabla_X \nabla_Y Z - \nabla_Y \nabla_X Z - \nabla_{[X,Y]} Z.
\]
In lokalen Koordinaten sind die Komponenten $R^l{}_{\!kij}$ definiert durch
$R(\partial_i, \partial_j)\partial_k = R^l{}_{\!kij}\,\partial_l$:
\[
  R^l{}_{\!kij} = \partial_i \Gamma^l_{jk} - \partial_j \Gamma^l_{ik}
  + \Gamma^l_{im}\Gamma^m_{jk} - \Gamma^l_{jm}\Gamma^m_{ik}.
\]
\end{definition}

\begin{proposition}[Symmetrien des Riemann-Tensors]
Mit herabgezogenen Indizes $R_{ijkl} = g_{im}R^m{}_{\!jkl}$:
\begin{enumerate}[label=(\roman*)]
  \item $R_{ijkl} = -R_{jikl}$ (Antisymmetrie im ersten Indexpaar),
  \item $R_{ijkl} = -R_{ijlk}$ (Antisymmetrie im zweiten Indexpaar),
  \item $R_{ijkl} = R_{klij}$ (Paarsymmetrie),
  \item $R_{ijkl} + R_{iklj} + R_{iljk} = 0$ (algebraische Bianchi-Identität).
\end{enumerate}
\end{proposition}

\subsection{Ricci- und Skalarkrümmung}

\begin{definition}
Der \emph{Ricci-Tensor} $\mathrm{Ric}$ ist die Kontraktion:
\[
  R_{ij} = R^k{}_{\!ikj} = g^{kl}R_{kilj}.
\]
Die \emph{Skalarkrümmung} ist die volle Spur:
\[
  R = g^{ij}R_{ij} = \sum_i R^i{}_i.
\]
\end{definition}

\begin{definition}[Schnittkrümmung]
Für eine 2-Ebene $\sigma = \mathrm{span}\{u, v\} \subset T_pM$ ist die
\emph{Schnittkrümmung}:
\[
  K(\sigma) = \frac{g(R(u,v)v, u)}{g(u,u)g(v,v) - g(u,v)^2}.
\]
\end{definition}

\begin{beispiel}[Krümmung klassischer Flächen]
\begin{enumerate}[label=(\alph*)]
  \item \textbf{Sphäre $S^n(r)$}: Konstante Schnittkrümmung $K = 1/r^2$.
  \item \textbf{Euklidischer Raum}: $K = 0$ überall.
  \item \textbf{Hyperbolischer Raum $\mathbb{H}^n$}: Konstant $K = -1$.
  \item \textbf{Flacher Torus $\mathbb{T}^2$}: $K = 0$.
\end{enumerate}
\end{beispiel}

\subsection{Gaußsche Krümmung über Fundamentalformen}

Für eine Fläche $\mathbf{r}(u,v)$ in $\mathbb{R}^3$ hat die \emph{zweite
Fundamentalform} die Koeffizienten:
\[
  L = \mathbf{r}_{uu} \cdot \mathbf{n}, \quad
  M = \mathbf{r}_{uv} \cdot \mathbf{n}, \quad
  N = \mathbf{r}_{vv} \cdot \mathbf{n},
\]
mit Einheitsnormale $\mathbf{n} = (\mathbf{r}_u \times \mathbf{r}_v)/|\mathbf{r}_u \times \mathbf{r}_v|$.
Die \emph{Gaußsche Krümmung} und \emph{mittlere Krümmung} sind:
\[
  K = \frac{LN - M^2}{EG - F^2}, \qquad
  H = \frac{EN + GL - 2FM}{2(EG - F^2)}.
\]

\begin{theorem}[Theorema Egregium --- Gauß, 1827]
Die Gaußsche Krümmung $K$ einer Fläche in $\mathbb{R}^3$ hängt nur von der
ersten Fundamentalform $E, F, G$ und deren Ableitungen ab. Insbesondere
ist $K$ invariant unter lokalen Isometrien.
\end{theorem}

\begin{bemerkung}
Dieses Theorem ist wahrhaftig bemerkenswert (wie sein lateinischer Name besagt):
Ein flaches Blatt Papier, zu einem Zylinder gerollt, hat sowohl flach als auch
zylindrisch $K = 0$. Dagegen hat $K \neq 0$ auf einer Sphäre zur Folge, dass
sie nicht isometrisch auf eine Ebene abgebildet werden kann --- daher verzerren
Landkarten zwangsläufig Flächen oder Winkel.
\end{bemerkung}

% ============================================================
\section{Der Gauß-Bonnet-Satz}
\label{sec:gauss-bonnet}
% ============================================================

Der Gauß-Bonnet-Satz ist eines der schönsten Resultate der Mathematik.
Er verknüpft den lokalen differentialgeometrischen Begriff der Krümmung mit
dem globalen topologischen Invarianten --- der Euler-Charakteristik.

\begin{theorem}[Gauß-Bonnet für geschlossene Flächen]
\label{thm:gauss-bonnet}
Sei $(M, g)$ eine kompakte orientierte Riemannsche 2-Mannigfaltigkeit ohne Rand.
Dann gilt:
\[
  \int_M K\,dA = 2\pi\,\chi(M),
\]
wobei $K$ die Gaußsche Krümmung, $dA$ das Flächenelement und $\chi(M)$
die Euler-Charakteristik von $M$ ist.
\end{theorem}

\begin{proof}[Beweisübersicht]
Wir verwenden ein kombinatorisch-analytisches Argument.

\textbf{Schritt 1 (Lokaler Gauß-Bonnet).} Für ein geodätisches Polygon $P$
mit Ecken $p_1, \ldots, p_k$, Innenwinkeln $\alpha_i$ und Außenwinkeln
$\theta_i = \pi - \alpha_i$ gilt:
\[
  \int_P K\,dA + \oint_{\partial P} \kappa_g\,ds = \sum_{i=1}^k (\pi - \alpha_i),
\]
wobei $\kappa_g$ die geodätische Krümmung von $\partial P$ ist.
Dies folgt aus der Gaußschen Abbildung und dem Satz von Stokes.

\textbf{Schritt 2 (Triangulierung).} Trianguliere $M$ durch geodätische Dreiecke
$T_1, \ldots, T_F$ mit $V$ Ecken und $E$ Kanten. Die lokale Formel wird auf
jedes Dreieck angewendet.

\textbf{Schritt 3 (Summation).} Summe über alle Dreiecke:
\[
  \int_M K\,dA = \sum_{i} \int_{T_i} K\,dA
  = (V - E + F) \cdot 2\pi = 2\pi\,\chi(M),
\]
unter Verwendung der Eulerschen Formel $\chi(M) = V - E + F$ und der
Tatsache, dass Innenwinkel um jeden Eckpunkt sich zu $2\pi$ summieren.
\end{proof}

\begin{korollar}[Topologische Konsequenzen]
\begin{enumerate}[label=(\roman*)]
  \item Für $S^2$: $\chi(S^2) = 2$, also $\int_{S^2} K\,dA = 4\pi$.
  \item Der Torus $\mathbb{T}^2$ hat $\chi = 0$, also $\int K\,dA = 0$:
        positive und negative Krümmungsbeiträge heben sich exakt auf.
  \item Keine Metrik auf $S^2$ kann $K \leq 0$ überall (mit strikter Ungleichung
        irgendwo) haben; keine Metrik auf $\mathbb{T}^2$ kann $K > 0$ haben.
\end{enumerate}
\end{korollar}

\subsection{Der Chern-Gauß-Bonnet-Satz}

\begin{theorem}[Chern-Gauß-Bonnet]
\label{thm:chern-gauss-bonnet}
Sei $M$ eine kompakte orientierte Riemannsche Mannigfaltigkeit gerader Dimension
$n = 2m$ ohne Rand. Dann gilt:
\[
  \chi(M) = \frac{1}{(2\pi)^m \, 2^m \, m!} \int_M \mathrm{Pf}(\Omega),
\]
wobei $\Omega$ die Krümmungs-2-Form des Levi-Civita-Zusammenhangs ist und
$\mathrm{Pf}$ das Pfaffian einer antisymmetrischen Matrix bezeichnet.
Für $n=4$ mit Weyl-Tensor $W$ und Ricci-Tensor $\mathrm{Ric}$:
\[
  \chi(M) = \frac{1}{8\pi^2}\int_M
  \left(|W|^2 - 2|\overset{\circ}{\mathrm{Ric}}|^2
  + \frac{R^2}{24}\right) dV,
\]
wobei $\overset{\circ}{\mathrm{Ric}}$ der spurfreie Anteil des Ricci-Tensors ist
(vgl. Besse, \emph{Einstein Manifolds}, Gl.~6.35).
\end{theorem}

Dies wurde von Shiing-Shen Chern 1944 mithilfe des Formalismus charakteristischer
Klassen und der Gauß-Bonnet-Chern-Form bewiesen, ohne eine Einbettung in einen
euklidischen Umgebungsraum vorauszusetzen.

% ============================================================
\section{Geodäten}
\label{sec:geodaeten}
% ============================================================

Geodäten sind die ``geraden Linien'' der Riemannschen Geometrie ---
Kurven, die lokal die Länge minimieren.

\begin{definition}[Geodäte]
Eine glatte Kurve $\gamma: I \to M$ heißt \emph{Geodäte}, falls ihr
Tangentialvektor parallel zu sich selbst ist:
\[
  \nabla_{\gamma'}\gamma' = 0.
\]
In lokalen Koordinaten, mit $\gamma(t) = (x^1(t), \ldots, x^n(t))$,
lautet dies die \emph{Geodätengleichung}:
\[
  \ddot{x}^k + \Gamma^k_{ij}\,\dot{x}^i\,\dot{x}^j = 0,
  \quad k = 1, \ldots, n.
\]
\end{definition}

\begin{proposition}
Geodäten sind Kurven konstanter Geschwindigkeit: $|\gamma'(t)|_g = \mathrm{const}$.
\end{proposition}

\begin{beispiel}[Geodäten auf klassischen Räumen]
\begin{enumerate}[label=(\alph*)]
  \item In $\mathbb{R}^n$: Geraden $\gamma(t) = p + tv$.
  \item Auf $S^2$: Großkreise (Schnitte von Ebenen durch den Ursprung).
  \item In $\mathbb{H}^2$: Senkrechte Geraden $x = \mathrm{const}$ und
        Halbkreise mit Durchmesser auf der $x$-Achse.
\end{enumerate}
\end{beispiel}

\subsection{Satz von Hopf-Rinow}

\begin{theorem}[Hopf-Rinow]
\label{thm:hopf-rinow}
Sei $(M, g)$ eine zusammenhängende Riemannsche Mannigfaltigkeit.
Die folgenden Aussagen sind äquivalent:
\begin{enumerate}[label=(\roman*)]
  \item $(M, d_g)$ ist ein vollständiger metrischer Raum (jede Cauchy-Folge konvergiert),
  \item $M$ ist \emph{geodätisch vollständig}: Für jeden Punkt $p \in M$ und jeden
        Vektor $v \in T_pM$ ist die Geodäte $\gamma_v$ für alle $t \in \mathbb{R}$
        definiert,
  \item Für einen (äquivalent: jeden) Punkt $p \in M$ ist die Exponentialabbildung
        $\exp_p$ auf ganz $T_pM$ definiert.
\end{enumerate}
Außerdem impliziert jede dieser Bedingungen:
\begin{enumerate}[label=(\alph*)]
  \item Je zwei Punkte in $M$ sind durch eine minimierende Geodäte verbunden.
\end{enumerate}
\end{theorem}

\begin{proof}[Beweisübersicht]
Die Implikation (i) $\Rightarrow$ (ii) verwendet, dass Geodäten in endlicher Zeit
keine ``Unendlichkeit erreichen'' können, wenn der metrische Raum vollständig ist:
Wäre eine Geodäte $\gamma$ nur auf $[0, T)$ definiert, dann wäre $(\gamma(t_n))$
für $t_n \nearrow T$ eine Cauchy-Folge, die gegen ein $p^* \in M$ konvergiert,
und die Geodäte lässt sich durch $p^*$ fortsetzen.

Die Implikation (iii) $\Rightarrow$ (a) zeigt man, indem man für ein beliebiges
$q \in M$ das Infimum des Längenfunktionals über alle Kurven von $p$ nach $q$
als angenommen nachweist: Man konstruiert eine minimierende Folge und extrahiert
mit der Geodätenkugelstruktur eine konvergente Teilfolge, dann nutzt man
$\exp_p$, das auf ganz $T_pM$ definiert ist, um die minimierende Geodäte
fortzusetzen.
\end{proof}

\subsection{Jacobi-Felder}

\begin{definition}[Jacobi-Feld]
Sei $\gamma$ eine Geodäte. Ein Vektorfeld $J$ längs $\gamma$ heißt
\emph{Jacobi-Feld}, wenn es die \emph{Jacobi-Gleichung} erfüllt:
\[
  \nabla^2_{\gamma'} J + R(\gamma', J)\gamma' = 0,
\]
also $J'' + R(\gamma', J)\gamma' = 0$ mit $J'' = \nabla_{\gamma'}\nabla_{\gamma'} J$.
\end{definition}

Jacobi-Felder beschreiben infinitesimale Geodäten-Variationen und kodieren,
wie benachbarte Geodäten konvergieren oder divergieren. Ein \emph{konjugierter
Punkt} zu $p$ längs $\gamma$ ist ein Punkt $\gamma(t_0)$, an dem ein
nicht-triviales Jacobi-Feld verschwindet. Auf $S^2$ ist der konjugierte
Ort des Nordpols der Südpol.

% ============================================================
\section{Untermannigfaltigkeiten}
\label{sec:untermannigfaltigkeiten}
% ============================================================

\begin{definition}[Eingebettete Untermannigfaltigkeit]
Eine \emph{eingebettete Untermannigfaltigkeit} von $(M, g)$ ist eine glatte
Untermannigfaltigkeit $\iota: \Sigma \hookrightarrow M$ mit der induzierten
(zurückgezogenen) Metrik $\bar{g} = \iota^* g$.
\end{definition}

\subsection{Die zweite Fundamentalform}

Für eine Hyperfläche $\Sigma \subset M$ mit Einheitsnormale $\mathbf{n}$
ist die \emph{zweite Fundamentalform} die symmetrische Bilinearform:
\[
  \mathrm{II}(X,Y) = g(\nabla_X Y, \mathbf{n}), \quad X, Y \in T\Sigma.
\]
Ihre Spur (bezüglich $\bar{g}$) ist die \emph{mittlere Krümmung}:
\[
  H = \frac{1}{n-1}\,\mathrm{tr}_{\bar{g}}(\mathrm{II}).
\]

\begin{definition}[Minimalfläche]
Eine Fläche $\Sigma$ heißt \emph{Minimalfläche}, wenn $H \equiv 0$.
\end{definition}

\begin{beispiel}[Klassische Minimalflächen]
\begin{enumerate}[label=(\alph*)]
  \item \textbf{Katenoid}: $\mathbf{r}(u,v) = (\cosh u \cos v, \cosh u \sin v, u)$.
        Die einzige Minimalfläche der Revolution.
  \item \textbf{Helikoid}: $\mathbf{r}(u,v) = (u\cos v, u\sin v, v)$.
        Einfach zusammenhängende Minimalfläche, isometrisch zum Katenoid.
  \item \textbf{Enneper-Fläche}: Dargestellt durch die Weierstraß-Enneper-Darstellung;
        selbstschneidend, aber minimal.
\end{enumerate}
\end{beispiel}

\subsection{Gauß- und Codazzi-Gleichungen}

Die Gauß-Gleichung verbindet die innere Krümmung von $\Sigma$ mit der Umgebungskrümmung
und der zweiten Fundamentalform:
\[
  \bar{R}(X,Y,Z,W) = R(X,Y,Z,W)
  + \mathrm{II}(X,W)\,\mathrm{II}(Y,Z) - \mathrm{II}(X,Z)\,\mathrm{II}(Y,W).
\]
Im klassischen Fall $M = \mathbb{R}^3$ mit $R = 0$ ergibt sich
$K = \kappa_1 \kappa_2$ (Produkt der Hauptkrümmungen) --- der Inhalt
des Theorema Egregium.

Die Codazzi-Mainardi-Gleichungen sind Integrabilitätsbedingungen:
\[
  (\nabla_X \mathrm{II})(Y,Z) = (\nabla_Y \mathrm{II})(X,Z) + R(X,Y)Z^{\perp}.
\]

% ============================================================
\section{Lie-Gruppen und Geometrie}
\label{sec:liegruppen}
% ============================================================

Lie-Gruppen vereinen die algebraische Struktur einer Gruppe mit der glatten
Struktur einer Mannigfaltigkeit und liefern eine reiche Quelle symmetrischer Räume.

\begin{definition}[Lie-Gruppe]
Eine \emph{Lie-Gruppe} ist eine Gruppe $G$, die gleichzeitig eine glatte
Mannigfaltigkeit ist, sodass die Gruppenoperationen $\mu: G \times G \to G$,
$(g,h) \mapsto gh$ und $\iota: G \to G$, $g \mapsto g^{-1}$ glatt sind.
\end{definition}

\begin{beispiel}
\begin{enumerate}[label=(\alph*)]
  \item $\mathbb{R}^n$ (Addition), $\mathbb{R}^* = \mathbb{R} \setminus \{0\}$,
  \item $\mathrm{GL}(n,\mathbb{R})$, $\mathrm{SL}(n,\mathbb{R})$, $\mathrm{O}(n)$, $\mathrm{SO}(n)$, $\mathrm{U}(n)$, $\mathrm{SU}(n)$,
  \item Die Kreisgruppe $S^1 \cong \mathrm{U}(1)$.
\end{enumerate}
\end{beispiel}

\begin{definition}[Lie-Algebra]
Die \emph{Lie-Algebra} $\mathfrak{g}$ von $G$ ist $T_eG$, ausgestattet mit
der Klammer $[X, Y] = \mathcal{L}_X Y$ (Lie-Ableitung). Für Matrix-Lie-Gruppen
besteht $\mathfrak{g}$ aus Matrizen mit $[A,B] = AB - BA$.
\end{definition}

\subsection{Linksinvariante Metriken und Killing-Felder}

Eine Riemannsche Metrik $g$ auf $G$ heißt \emph{linksinvariant}, wenn
$L_h^* g = g$ für alle $h \in G$, mit $L_h(x) = hx$. Eine solche Metrik
ist eindeutig durch ihren Wert bei $e$ bestimmt, also durch ein Skalarprodukt
auf $\mathfrak{g}$.

\begin{definition}[Killing-Vektorfeld]
Ein Vektorfeld $X$ auf $(M, g)$ heißt \emph{Killing-Feld}, wenn
$\mathcal{L}_X g = 0$, d.\,h.\ $X$ erzeugt Isometrien. Äquivalent:
\[
  \nabla_i X_j + \nabla_j X_i = 0.
\]
\end{definition}

\begin{proposition}
Auf einer Lie-Gruppe $G$ mit linksinvarianter Metrik sind die linksinvarianten
Vektorfelder genau dann Killing-Felder, wenn die Metrik auch rechtsinvariant
(\emph{bi-invariant}) ist.
\end{proposition}

\begin{beispiel}
Die bi-invariante Metrik auf $\mathrm{SO}(3)$ ist $\langle A, B \rangle = -\tfrac{1}{2}\mathrm{tr}(AB)$.
Geodäten sind Einparameter-Untergruppen $t \mapsto e^{tX}$, und die
Schnittkrümmung erfüllt $K(\sigma) \geq 0$.
\end{beispiel}

% ============================================================
\section{Anwendungen in der Allgemeinen Relativitätstheorie}
\label{sec:anwendungen}
% ============================================================

Albert Einsteins Allgemeine Relativitätstheorie (1915) ist wohl die bedeutendste
Anwendung der Differentialgeometrie in der Physik. Raumzeit wird als
4-dimensionale \emph{Lorentz-Mannigfaltigkeit} $(M, g)$ mit Metriksignatur
$(-,+,+,+)$ modelliert, und Gravitation wird mit Raumzeitkrümmung identifiziert.

\subsection{Lorentz-Mannigfaltigkeiten}

\begin{definition}
Eine \emph{Lorentz-Mannigfaltigkeit} $(M, g)$ ist eine glatte Mannigfaltigkeit
mit einem metrischen Tensor $g$ der Signatur $(-,+,+,+)$. Vektoren
$v \in T_pM$ werden eingeteilt in:
\begin{itemize}
  \item \emph{zeitartig}: $g(v,v) < 0$ (massive Teilchen),
  \item \emph{lichtartig (null)}: $g(v,v) = 0$ (Photonen),
  \item \emph{raumartig}: $g(v,v) > 0$.
\end{itemize}
\end{definition}

\subsection{Einsteinsche Feldgleichungen}

\begin{definition}[Einstein-Tensor]
Der \emph{Einstein-Tensor} ist:
\[
  G_{\mu\nu} = R_{\mu\nu} - \tfrac{1}{2} R\,g_{\mu\nu},
\]
wobei $R_{\mu\nu}$ der Ricci-Tensor und $R$ die Skalarkrümmung ist.
Er erfüllt die kontrahierte Bianchi-Identität $\nabla^\mu G_{\mu\nu} = 0$.
\end{definition}

\begin{theorem}[Einsteinsche Feldgleichungen]
In der Allgemeinen Relativitätstheorie wird die Geometrie der Raumzeit durch
die Verteilung von Materie und Energie gemäß den Gleichungen bestimmt:
\[
  G_{\mu\nu} + \Lambda g_{\mu\nu} = \frac{8\pi G}{c^4}\,T_{\mu\nu},
\]
wobei $\Lambda$ die kosmologische Konstante, $G$ die Newtonsche Gravitationskonstante
und $T_{\mu\nu}$ der Energie-Impuls-Tensor ist. Im Vakuum ($T_{\mu\nu} = 0$,
$\Lambda = 0$) vereinfacht sich dies zu $R_{\mu\nu} = 0$ (\emph{Ricci-flach}).
\end{theorem}

\subsection{Geodäten als Weltlinien}

\begin{proposition}
In der Allgemeinen Relativitätstheorie bewegen sich frei fallende Teilchen
(die keiner nicht-gravitativen Kraft ausgesetzt sind) auf \emph{zeitartigen
Geodäten} der Raumzeitmetrik. Photonen folgen \emph{lichtartigen Geodäten}.
\end{proposition}

\begin{beispiel}[Schwarzschild-Geodäten]
Die Schwarzschild-Metrik (außerhalb einer kugelsymmetrischen Masse $m$) lautet:
\[
  ds^2 = -\left(1 - \frac{r_s}{r}\right)c^2\,dt^2
  + \left(1 - \frac{r_s}{r}\right)^{-1}dr^2
  + r^2(d\theta^2 + \sin^2\!\theta\,d\phi^2),
\]
mit Schwarzschild-Radius $r_s = 2Gm/c^2$. Kreisförmige Geodäten existieren
bei $r = 3r_s/2$ (Photonensphäre) und $r = 3r_s$ (innerste stabile Kreisbahn).
Die Periheldrehung des Merkur ($43''$ pro Jahrhundert) ist eine direkte
Vorhersage aus den Geodätengleichungen.
\end{beispiel}

\subsection{Verbindung zum Python-Modul}

Das begleitende Modul \texttt{differential\_geometry.py} implementiert:
\begin{itemize}
  \item \texttt{ParametricCurve}: Frenet-Serret-Rahmen, Krümmung $\kappa$, Torsion $\tau$,
  \item \texttt{ParametricSurface}: erste und zweite Fundamentalform, $K$, $H$,
  \item \texttt{RiemannianGeometry}: Christoffel-Symbole $\Gamma^k_{ij}$,
        Riemann-Tensor $R^l{}_{kij}$, Ricci-Tensor $R_{ij}$, Skalarkrümmung,
  \item \texttt{GeodesicComputation}: numerische Geodäten-Integration via RK4/RK45,
  \item \texttt{ClassicalSurfaces}: Sphäre, Torus, Katenoid, Helikoid, Sattelfläche.
\end{itemize}
Krümmungsberechnungen verwenden die Koszul-Formel für Christoffel-Symbole und
kontrahieren entsprechend. Geodäten werden als gewöhnliche Differentialgleichungen
zweiter Ordnung mit \texttt{scipy.integrate.solve\_ivp} (Dormand-Prince-Methode)
gelöst.

% ============================================================
\section*{Literatur}
% ============================================================

\begin{thebibliography}{99}

\bibitem{docarmo1976}
M.~P. do~Carmo,
\textit{Differentialgeometrie von Kurven und Flächen},
Vieweg, 1983. (Originalausgabe: Differential Geometry of Curves and Surfaces,
Prentice-Hall, 1976.)

\bibitem{docarmo1992}
M.~P. do~Carmo,
\textit{Riemannian Geometry},
Birkh\"auser, Boston, 1992.

\bibitem{spivak1970}
M.~Spivak,
\textit{A Comprehensive Introduction to Differential Geometry},
Bd.~1--5, Publish or Perish, 1970--1975.

\bibitem{lee1997}
J.~M. Lee,
\textit{Riemannian Manifolds: An Introduction to Curvature},
Springer, New York, 1997.

\bibitem{lee2013}
J.~M. Lee,
\textit{Introduction to Smooth Manifolds},
2.~Aufl., Springer, New York, 2013.

\bibitem{oneill1983}
B.~O'Neill,
\textit{Semi-Riemannian Geometry with Applications to Relativity},
Academic Press, 1983.

\bibitem{petersen2006}
P.~Petersen,
\textit{Riemannian Geometry},
2.~Aufl., Springer, New York, 2006.

\bibitem{chern1944}
S.-S. Chern,
A simple intrinsic proof of the Gauss-Bonnet formula for closed Riemannian manifolds,
\textit{Ann.\ of Math.}, \textbf{45} (1944), 747--752.

\bibitem{gallot2004}
S.~Gallot, D.~Hulin, J.~Lafontaine,
\textit{Riemannian Geometry},
3.~Aufl., Springer, Berlin, 2004.

\bibitem{hopfrinow1931}
H.~Hopf, W.~Rinow,
Über den Begriff der vollständigen differentialgeometrischen Fläche,
\textit{Comment.\ Math.\ Helv.}, \textbf{3} (1931), 209--225.

\end{thebibliography}

\end{document}
